% Options for packages loaded elsewhere
\PassOptionsToPackage{unicode}{hyperref}
\PassOptionsToPackage{hyphens}{url}
\documentclass[
]{article}
\usepackage{xcolor}
\usepackage{amsmath,amssymb}
\setcounter{secnumdepth}{-\maxdimen} % remove section numbering
\usepackage{iftex}
\ifPDFTeX
  \usepackage[T1]{fontenc}
  \usepackage[utf8]{inputenc}
  \usepackage{textcomp} % provide euro and other symbols
\else % if luatex or xetex
  \usepackage{unicode-math} % this also loads fontspec
  \defaultfontfeatures{Scale=MatchLowercase}
  \defaultfontfeatures[\rmfamily]{Ligatures=TeX,Scale=1}
\fi
\usepackage{lmodern}
\ifPDFTeX\else
  % xetex/luatex font selection
\fi
% Use upquote if available, for straight quotes in verbatim environments
\IfFileExists{upquote.sty}{\usepackage{upquote}}{}
\IfFileExists{microtype.sty}{% use microtype if available
  \usepackage[]{microtype}
  \UseMicrotypeSet[protrusion]{basicmath} % disable protrusion for tt fonts
}{}
\makeatletter
\@ifundefined{KOMAClassName}{% if non-KOMA class
  \IfFileExists{parskip.sty}{%
    \usepackage{parskip}
  }{% else
    \setlength{\parindent}{0pt}
    \setlength{\parskip}{6pt plus 2pt minus 1pt}}
}{% if KOMA class
  \KOMAoptions{parskip=half}}
\makeatother
\usepackage{longtable,booktabs,array}
\usepackage{calc} % for calculating minipage widths
% Correct order of tables after \paragraph or \subparagraph
\usepackage{etoolbox}
\makeatletter
\patchcmd\longtable{\par}{\if@noskipsec\mbox{}\fi\par}{}{}
\makeatother
% Allow footnotes in longtable head/foot
\IfFileExists{footnotehyper.sty}{\usepackage{footnotehyper}}{\usepackage{footnote}}
\makesavenoteenv{longtable}
\setlength{\emergencystretch}{3em} % prevent overfull lines
\providecommand{\tightlist}{%
  \setlength{\itemsep}{0pt}\setlength{\parskip}{0pt}}
\usepackage{bookmark}
\IfFileExists{xurl.sty}{\usepackage{xurl}}{} % add URL line breaks if available
\urlstyle{same}
\hypersetup{
  hidelinks,
  pdfcreator={LaTeX via pandoc}}

\author{}
\date{}

\begin{document}

{
\setcounter{tocdepth}{3}
\tableofcontents
}
\section{ENVIRO\_DL STM32 --- User
Guide}\label{enviro_dl-stm32-user-guide}

\textbf{Product:} ENVIRO\_DL Environmental Data Logger
\textbf{Platform:} STM32F417 (ARM Cortex-M4) \textbf{Firmware:}
ENVIRO\_DL\_STM32F \textbf{Manufacturer:} EIPL (Envirotech Instruments
Pvt. Ltd.)

\begin{center}\rule{0.5\linewidth}{0.5pt}\end{center}

\subsection{Table of Contents}\label{table-of-contents}

\begin{enumerate}
\def\labelenumi{\arabic{enumi}.}
\tightlist
\item
  \hyperref[1-system-overview]{System Overview}
\item
  \hyperref[2-hardware-description]{Hardware Description}
\item
  \hyperref[3-getting-started]{Getting Started}
\item
  \hyperref[4-front-panel-operation]{Front Panel Operation}
\item
  \hyperref[5-lcd-display-pages]{LCD Display Pages}
\item
  \hyperref[6-sensor-suite]{Sensor Suite}
\item
  \hyperref[7-data-logging]{Data Logging}
\item
  \hyperref[8-configuration]{Configuration}
\item
  \hyperref[9-communication-interfaces]{Communication Interfaces}
\item
  \hyperref[10-terminal-command-reference]{Terminal Command Reference}
\item
  \hyperref[11-remote-data-upload]{Remote Data Upload}
\item
  \hyperref[12-firmware-over-the-air-fota-update]{Firmware Over-the-Air
  (FOTA) Update}
\item
  \hyperref[13-led-status-indicator]{LED Status Indicator}
\item
  \hyperref[14-fault-codes--troubleshooting]{Fault Codes \&
  Troubleshooting}
\item
  \hyperref[15-sd-card-file-reference]{SD Card File Reference}
\item
  \hyperref[16-technical-specifications]{Technical Specifications}
\item
  \hyperref[17-known-limitations]{Known Limitations}
\end{enumerate}

\begin{center}\rule{0.5\linewidth}{0.5pt}\end{center}

\subsection{1. System Overview}\label{system-overview}

The \textbf{ENVIRO\_DL} is a rugged, embedded environmental data logger
designed for continuous outdoor monitoring. It measures a comprehensive
set of meteorological parameters, stores readings locally on an SD card,
and transmits data to a remote server via an integrated LTE modem.

\subsubsection{Key Capabilities}\label{key-capabilities}

\begin{itemize}
\tightlist
\item
  Measures temperature, humidity, pressure, wind speed/direction,
  rainfall, and solar radiation
\item
  Logs data to an SD card in daily CSV files at a configurable interval
  (default: every 5 minutes)
\item
  Uploads data to a remote server over LTE (HTTP POST), configurable
  upload period
\item
  Displays live readings and system status on a 128×64 graphical LCD
\item
  Accepts configuration and commands via USB (virtual COM port) or UART
  terminal
\item
  Supports Firmware Over-the-Air (FOTA) updates with automatic rollback
  on failure
\item
  Maintains time across power loss using an external battery-backed RTC
  (MCP7940)
\item
  Runs FreeRTOS with multiple concurrent tasks for reliable, real-time
  operation
\end{itemize}

\begin{center}\rule{0.5\linewidth}{0.5pt}\end{center}

\subsection{2. Hardware Description}\label{hardware-description}

\subsubsection{Block Diagram}\label{block-diagram}

\begin{verbatim}
 ┌──────────────────────────────────────────────────────────────────┐
 │                        STM32F417 MCU                            │
 │                                                                  │
 │  ┌──────────┐  ┌──────────┐  ┌──────────┐  ┌───────────────┐   │
 │  │SensorTask│  │UploadTask│  │ CmdTask  │  │ KeyboardTask  │   │
 │  └────┬─────┘  └────┬─────┘  └────┬─────┘  └──────┬────────┘   │
 │       │              │              │                │            │
 └───────┼──────────────┼──────────────┼────────────────┼───────────┘
         │              │              │                │
    ┌────▼────┐    ┌────▼────┐   ┌────▼────┐   ┌───────▼───────┐
    │ Sensors │    │LTE Modem│   │USB / UART│  │Keypad + Pulses│
    │I2C/UART │    │(Quectel)│   │(Terminal)│  │Front Panel    │
    └─────────┘    └─────────┘   └──────────┘  └───────────────┘
         │
    ┌────▼──────────────────────────────────────────────────┐
    │  NAU7802 (Temp) | ADS112C04 (Hum/P/Solar/WindDir)    │
    │  RS485 Wind Sensor | SHT25/45/BMP280 | Rain Gauge    │
    └───────────────────────────────────────────────────────┘
\end{verbatim}

\subsubsection{Connectors \& Ports}\label{connectors-ports}

\begin{longtable}[]{@{}
  >{\raggedright\arraybackslash}p{(\linewidth - 4\tabcolsep) * \real{0.2727}}
  >{\raggedright\arraybackslash}p{(\linewidth - 4\tabcolsep) * \real{0.2727}}
  >{\raggedright\arraybackslash}p{(\linewidth - 4\tabcolsep) * \real{0.4545}}@{}}
\toprule\noalign{}
\begin{minipage}[b]{\linewidth}\raggedright
Port
\end{minipage} & \begin{minipage}[b]{\linewidth}\raggedright
Type
\end{minipage} & \begin{minipage}[b]{\linewidth}\raggedright
Function
\end{minipage} \\
\midrule\noalign{}
\endhead
\bottomrule\noalign{}
\endlastfoot
USB & USB-B / Micro-USB & Virtual COM port (CDC) for terminal commands
\& debug \\
UART1 & RS-232 / TTL header & Quectel LTE modem; also alternate terminal
interface \\
UART2 / RS485 & 2-wire RS485 & Wind speed sensor (Modbus RTU or Gill
ASCII) \\
I2C (Internal) & PCB traces & NAU7802, ADS112C04, SHT25/45, BMP280,
MCP7940 RTC \\
SD Card & MicroSD slot & Data logging, configuration, FOTA firmware
staging \\
Rain Gauge & 2-pin screw terminal & Tipping bucket pulse input (PG7) \\
LCD & On-board ribbon & 128×64 GLCD display \\
Keypad & On-board or header & 6-button navigation
(Up/Down/Left/Right/ESC/Enter) \\
SIM Card & SIM slot (modem) & LTE SIM for data connectivity \\
LED & On-board & Green/Red status indicator (PG4) \\
\end{longtable}

\begin{center}\rule{0.5\linewidth}{0.5pt}\end{center}

\subsection{3. Getting Started}\label{getting-started}

\subsubsection{3.1 Prerequisites}\label{prerequisites}

Before powering on, ensure:

\begin{enumerate}
\def\labelenumi{\arabic{enumi}.}
\tightlist
\item
  \textbf{SD card inserted} --- FAT32 formatted, minimum 1 GB
  recommended
\item
  \textbf{SIM card installed} in the LTE modem slot (if remote upload is
  required)
\item
  \textbf{Sensors connected} per the wiring diagram (rain gauge, wind
  sensor, analog sensors)
\item
  \textbf{Antenna attached} to the LTE modem SMA connector
\end{enumerate}

\subsubsection{3.2 First Power-On}\label{first-power-on}

\begin{enumerate}
\def\labelenumi{\arabic{enumi}.}
\tightlist
\item
  Apply 5--12 V DC power to the unit
\item
  The LCD will display the \textbf{Envirotech splash screen} for
  approximately 5 seconds
\item
  The system will:

  \begin{itemize}
  \tightlist
  \item
    Synchronize time from the external RTC (MCP7940) if previously set
  \item
    Mount the SD card and load \texttt{config.bin}, \texttt{nvdata.bin}
  \item
    Initialize all sensors and verify communication
  \item
    Start the FreeRTOS tasks (sensor reading, upload, UI, command)
  \end{itemize}
\item
  The STATUS LED will begin \textbf{blinking at 1 Hz} to indicate
  healthy operation
\item
  The LCD will transition to the \textbf{status dashboard} (Page 0)
\end{enumerate}

\subsubsection{3.3 Initial Configuration
(Required)}\label{initial-configuration-required}

Connect via USB (see \hyperref[91-usb-cdc-virtual-com-port]{Section
9.1}). All writable configuration is done through the \textbf{C02}
system config command. Changes are saved to the SD card automatically
when the command is accepted.

Set station identity:

\begin{verbatim}
(?eipl,C02,000000,sid=AWS_ROOF&lkey=SITE01?)
\end{verbatim}

Set the primary server IP:

\begin{verbatim}
(?eipl,C02,000000,sip1=192.168.1.100?)
\end{verbatim}

Set the current date and time (format: YYMMDDHHmm, 10 digits):

\begin{verbatim}
(?eipl,C02,000000,rtc=2603091030?)
\end{verbatim}

Read back station info to confirm:

\begin{verbatim}
(?eipl,C01,000000,?)
\end{verbatim}

After configuration, the unit will begin logging and uploading
automatically.

\begin{center}\rule{0.5\linewidth}{0.5pt}\end{center}

\subsection{4. Front Panel Operation}\label{front-panel-operation}

\subsubsection{4.1 Button Functions}\label{button-functions}

\begin{longtable}[]{@{}lll@{}}
\toprule\noalign{}
Button & Short Press & Long Press \\
\midrule\noalign{}
\endhead
\bottomrule\noalign{}
\endlastfoot
\textbf{Up} & Previous LCD page & --- \\
\textbf{Down} & Next LCD page & --- \\
\textbf{Left} & Previous item / Decrease value & --- \\
\textbf{Right} & Next item / Increase value & --- \\
\textbf{Enter} & Confirm selection & Enter configuration menu \\
\textbf{ESC} & Cancel / Exit menu & Return to main dashboard \\
\end{longtable}

\subsubsection{4.2 Navigation Flow}\label{navigation-flow}

\begin{verbatim}
Page 0: Status Bar
   ↕ (Up/Down)
Page 1: Temperature & Humidity
   ↕
Page 2: Wind Speed & Direction
   ↕
Page 3: Rainfall & Solar Radiation
   ↕
Page 4: Uplink Counters
   ↕
Page 5: FOTA Progress (only during update)
\end{verbatim}

\begin{center}\rule{0.5\linewidth}{0.5pt}\end{center}

\subsection{5. LCD Display Pages}\label{lcd-display-pages}

\subsubsection{Page 0 --- System Status
Bar}\label{page-0-system-status-bar}

\begin{verbatim}
 [ANT] GSM:▌▌▌  Bat:12.3V
 2026-03-09  10:25:00
 SD:OK  MOD:OK  RTC:OK
\end{verbatim}

\begin{longtable}[]{@{}ll@{}}
\toprule\noalign{}
Symbol & Meaning \\
\midrule\noalign{}
\endhead
\bottomrule\noalign{}
\endlastfoot
\texttt{{[}ANT{]}} & Antenna / LTE connected \\
\texttt{GSM:▌▌▌} & Signal strength (1--5 bars) \\
\texttt{Bat:12.3V} & Battery/supply voltage \\
\texttt{SD:OK} & SD card healthy \\
\texttt{MOD:OK} & LTE modem responding \\
\texttt{RTC:OK} & Real-time clock healthy \\
\end{longtable}

\subsubsection{Page 1 --- Temperature \&
Humidity}\label{page-1-temperature-humidity}

\begin{verbatim}
 Temp:  25.3 °C
 Hum:   64.5 %RH
 DewPt: 17.9 °C
 Press: 1013.2 mbar
\end{verbatim}

\subsubsection{Page 2 --- Wind}\label{page-2-wind}

\begin{verbatim}
 Wind Spd:  3.2 m/s
 Wind Dir:  NNE (025°)
 Gust:      5.1 m/s
 Avg(5min): 2.9 m/s
\end{verbatim}

\subsubsection{Page 3 --- Rainfall \&
Solar}\label{page-3-rainfall-solar}

\begin{verbatim}
 Rain Today: 12.50 mm
 Rain/Hour:   2.00 mm
 Solar:      450 W/m²
\end{verbatim}

\subsubsection{Page 4 --- Upload Counters}\label{page-4-upload-counters}

\begin{verbatim}
 Uploads Attempted: 145
 Uploads Success:   143
 Last Upload: 10:20:00
 Next Upload: 10:25:00
\end{verbatim}

\subsubsection{Page 5 --- FOTA Progress (during update
only)}\label{page-5-fota-progress-during-update-only}

\begin{verbatim}
 FOTA Update in Progress
 [##########          ] 52%
 Downloading firmware...
 DO NOT POWER OFF
\end{verbatim}

\begin{center}\rule{0.5\linewidth}{0.5pt}\end{center}

\subsection{6. Sensor Suite}\label{sensor-suite}

\subsubsection{6.1 Temperature --- NAU7802 + PT1000
RTD}\label{temperature-nau7802-pt1000-rtd}

\begin{itemize}
\tightlist
\item
  \textbf{Type:} PT1000 resistance temperature detector (RTD)
\item
  \textbf{Interface:} I2C, address 0x54
\item
  \textbf{Resolution:} 24-bit ADC
\item
  \textbf{Accuracy:} ±0.3 °C (sensor dependent)
\item
  \textbf{Conversion:} Callendar-Van Dusen polynomial
\item
  \textbf{Output channels:} ch2 (primary temperature)
\end{itemize}

\subsubsection{6.2 Humidity, Pressure, Solar Radiation, Wind Direction
---
ADS112C04}\label{humidity-pressure-solar-radiation-wind-direction-ads112c04}

\begin{itemize}
\tightlist
\item
  \textbf{Type:} 24-bit multi-channel ADC (Texas Instruments ADS112C04)
\item
  \textbf{Interface:} I2C
\item
  \textbf{Channels used:}
\end{itemize}

\begin{longtable}[]{@{}lll@{}}
\toprule\noalign{}
ADC Channel & Parameter & Unit \\
\midrule\noalign{}
\endhead
\bottomrule\noalign{}
\endlastfoot
Ch0 & Relative Humidity & \%RH \\
Ch1 & Atmospheric Pressure & mbar \\
Ch3 & Solar Radiation & W/m² \\
Ch4 & Wind Direction (voltage) & degrees \\
\end{longtable}

\begin{itemize}
\tightlist
\item
  \textbf{Reference voltage:} Internal 2.048 V or external
\item
  \textbf{Calibration:} Gain and offset configurable via C02 command
\end{itemize}

\subsubsection{6.3 Wind Speed --- RS485
Sensor}\label{wind-speed-rs485-sensor}

\begin{itemize}
\tightlist
\item
  \textbf{Interface:} UART2 via RS485 transceiver
\item
  \textbf{Supported protocols:}
\end{itemize}

\begin{longtable}[]{@{}llll@{}}
\toprule\noalign{}
Protocol & Sensor Brand & Baud Rate & Format \\
\midrule\noalign{}
\endhead
\bottomrule\noalign{}
\endlastfoot
Modbus RTU & RENKEE & 9600 8N1 & Register query \\
ASCII (Gill) & Gill Instruments & Configurable &
\texttt{"w?\textbackslash{}r\textbackslash{}n"} query \\
\end{longtable}

\begin{itemize}
\tightlist
\item
  \textbf{Sample rate:} Approximately once per second
\item
  \textbf{Computed values:} Minimum, maximum, and average wind speed
  over the snap interval; gust speed
\item
  \textbf{Squall detection:} WMO-standard state machine; triggers on
  \textgreater25\% speed increase
\end{itemize}

\subsubsection{6.4 Rainfall --- Tipping Bucket
Gauge}\label{rainfall-tipping-bucket-gauge}

\begin{itemize}
\tightlist
\item
  \textbf{Interface:} GPIO falling-edge pulse on pin PG7
\item
  \textbf{Debounce:} Software debounce --- pulse is confirmed by
  re-reading the pin 20 ms later
\item
  \textbf{Default resolution:} 0.5 mm per tip (configurable via
  \texttt{rain\_mm\_per\_tick})
\item
  \textbf{Accumulators:}

  \begin{itemize}
  \tightlist
  \item
    \texttt{rain\_inday} --- resets at midnight (or configured reset
    hour)
  \item
    \texttt{rain\_inhour} --- resets after each upload period
  \end{itemize}
\item
  \textbf{Persistence:} Rain counts are saved to SD card
  (\texttt{nvdata.bin}) and survive power loss
\end{itemize}

\subsubsection{6.5 Secondary T/H/P --- SHT25 / SHT45 /
BMP280}\label{secondary-thp-sht25-sht45-bmp280}

\begin{itemize}
\tightlist
\item
  \textbf{Purpose:} Secondary/redundant temperature, humidity, and
  pressure measurement
\item
  \textbf{Interface:} I2C

  \begin{itemize}
  \tightlist
  \item
    SHT25: 0x40
  \item
    SHT45: 0x44
  \item
    BMP280: 0x76
  \end{itemize}
\item
  \textbf{Sensor type selected at firmware compile time}
\item
  \textbf{User-configurable offsets} applied to all readings (see C09
  command)
\end{itemize}

\subsubsection{6.6 Dew Point (Computed)}\label{dew-point-computed}

Dew point is computed automatically from temperature and relative
humidity using the Magnus formula. It is not a separate sensor and is
not currently included as a named channel in the upload packet or CSV
log. It is displayed locally on the LCD.

\begin{center}\rule{0.5\linewidth}{0.5pt}\end{center}

\subsection{7. Data Logging}\label{data-logging}

\subsubsection{7.1 Log File Format}\label{log-file-format}

Data is written to the SD card as daily \textbf{CSV files} in URL
query-string format:

\textbf{Filename:} \texttt{DDMMYY.csv} Example: \texttt{090326.csv} = 9
March 2026

\textbf{Each log line (one per snap interval):}

\begin{verbatim}
000001&locationkey=SITE01&stationkey=AWS_ROOF&datetime=1741513500&ch1=26.50&ch2=25.30&ch3=24.10&ch4=65.20&ch5=64.50&ch6=63.80&ch7=450.0&ch8=1013.2&ch9=025&ch10=4.80&ch11=3.20&ch12=2.50&ch13=12.50&ch14=12.3
\end{verbatim}

\subsubsection{7.2 Channel Reference}\label{channel-reference}

The upload packet and CSV log use the following fixed channel
assignments:

\begin{longtable}[]{@{}
  >{\raggedright\arraybackslash}p{(\linewidth - 6\tabcolsep) * \real{0.2647}}
  >{\raggedright\arraybackslash}p{(\linewidth - 6\tabcolsep) * \real{0.3235}}
  >{\raggedright\arraybackslash}p{(\linewidth - 6\tabcolsep) * \real{0.1765}}
  >{\raggedright\arraybackslash}p{(\linewidth - 6\tabcolsep) * \real{0.2353}}@{}}
\toprule\noalign{}
\begin{minipage}[b]{\linewidth}\raggedright
Channel
\end{minipage} & \begin{minipage}[b]{\linewidth}\raggedright
Parameter
\end{minipage} & \begin{minipage}[b]{\linewidth}\raggedright
Unit
\end{minipage} & \begin{minipage}[b]{\linewidth}\raggedright
Source
\end{minipage} \\
\midrule\noalign{}
\endhead
\bottomrule\noalign{}
\endlastfoot
ch1 & Temperature Maximum (AT\_Max) & °C* & PT1000/NAU7802, MMA over
snap interval \\
ch2 & Temperature Average (AT\_Avg) & °C* & PT1000/NAU7802, MMA over
snap interval \\
ch3 & Temperature Minimum (AT\_Min) & °C* & PT1000/NAU7802, MMA over
snap interval \\
ch4 & Humidity Maximum (RH\_Max) & \%RH & ADS112C04, MMA over snap
interval \\
ch5 & Humidity Average (RH\_Avg) & \%RH & ADS112C04, MMA over snap
interval \\
ch6 & Humidity Minimum (RH\_Min) & \%RH & ADS112C04, MMA over snap
interval \\
ch7 & Solar Radiation (instantaneous) & W/m²* & ADS112C04, last reading
at snap \\
ch8 & Barometric Pressure & mbar* & ADS112C04 / BMP280, last reading at
snap \\
ch9 & Wind Direction & degrees & RS485 sensor, last reading at snap \\
ch10 & Wind Speed Maximum (WS\_Max) & m/s* & RS485 sensor, MMA over snap
interval \\
ch11 & Wind Speed Average (WS\_Avg) & m/s* & RS485 sensor, MMA over snap
interval \\
ch12 & Wind Speed Minimum (WS\_Min) & m/s* & RS485 sensor, MMA over snap
interval \\
ch13 & Daily Rainfall & mm & Tipping bucket, accumulated since
midnight \\
ch14 & Battery / Supply Voltage & V & Always included \\
\end{longtable}

\begin{quote}
* Unit conversion applies if configured. See
\hyperref[84-unit-conversion]{Section 8.4}. Channels for disabled
sensors are omitted from the packet (not sent as zero).
\end{quote}

\subsubsection{7.3 Snap Interval}\label{snap-interval}

The \textbf{snap interval} defines how often a line is written to the
log file and how often data is uploaded. Default: \textbf{5 minutes}.

\begin{itemize}
\tightlist
\item
  Configurable from 1 to 1440 minutes
\item
  At each snap, the system computes min/max/average over all readings
  collected since the last snap
\item
  Rain accumulator reflects all tips since the start of the current day
\end{itemize}

\subsubsection{7.4 Min/Max/Average
Computation}\label{minmaxaverage-computation}

Within each snap interval the firmware continuously accumulates: -
Temperature, humidity, pressure, solar: simple running average - Wind
speed: separate min, max, average (1 sample per second) - Wind gust:
peak instantaneous speed, with timestamp and direction

All accumulators reset at the start of each new snap window.

\begin{center}\rule{0.5\linewidth}{0.5pt}\end{center}

\subsection{8. Configuration}\label{configuration}

Configuration is stored in binary files on the SD card and loaded at
boot. Changes made via C02 and C03 are \textbf{saved to SD card
automatically} --- no separate save command is needed.

\subsubsection{8.1 Station Info --- Read Only
(C01)}\label{station-info-read-only-c01}

C01 is a \textbf{read-only} status report. Send it without parameters to
retrieve the current station ID, firmware version, signal level, and
fault code:

\begin{verbatim}
(?eipl,C01,000000,?)
\end{verbatim}

\textbf{Response fields:}

\begin{longtable}[]{@{}ll@{}}
\toprule\noalign{}
Field & Description \\
\midrule\noalign{}
\endhead
\bottomrule\noalign{}
\endlastfoot
\texttt{SID} & Station/site ID \\
\texttt{T} & Current date-time (YYMMDDHHmm) \\
\texttt{DID} & Device ID \\
\texttt{FW} & Firmware version string \\
\texttt{Q} & GSM signal / BER \\
\texttt{F} & Fault bitmask (see Section 14.1) \\
\end{longtable}

\subsubsection{8.2 System Configuration --- Read / Write
(C02)}\label{system-configuration-read-write-c02}

C02 is the main configuration command. Send one or more
\texttt{key=value} pairs in a single packet. Changes are saved
automatically.

\begin{longtable}[]{@{}lll@{}}
\toprule\noalign{}
Parameter & Description & Example \\
\midrule\noalign{}
\endhead
\bottomrule\noalign{}
\endlastfoot
\texttt{sid} & Station/site ID (3--9 chars) & \texttt{sid=AWS\_ROOF} \\
\texttt{lkey} & Location key (1--6 chars) & \texttt{lkey=SITE01} \\
\texttt{mpm} & Snap interval in minutes (1--1440) & \texttt{mpm=5} \\
\texttt{upm} & Upload period in minutes (1--1440) & \texttt{upm=5} \\
\texttt{sip1} & Primary server IP address &
\texttt{sip1=192.168.1.100} \\
\texttt{sip2} & Secondary server IP address &
\texttt{sip2=192.168.1.101} \\
\texttt{rtc} & Set RTC date/time (YYMMDDHHmm, 10 digits) &
\texttt{rtc=2603091030} \\
\texttt{drh} & Daily rain reset hour (0--2350, HHMM) & \texttt{drh=0} \\
\texttt{rt} & Rain gauge tick value in mm (e.g.~0.5) &
\texttt{rt=0.5} \\
\texttt{smn} & Server mobile number (10--13 digits) &
\texttt{smn=9440520222} \\
\texttt{dmn1} & Device mobile number 1 & \texttt{dmn1=9876543210} \\
\texttt{dmn2} & Device mobile number 2 & \texttt{dmn2=9123456789} \\
\texttt{ucf} & Uplink control flags (hex byte) & \texttt{ucf=80} (GPRS
ON) \\
\end{longtable}

\textbf{Examples:}

\begin{verbatim}
# Set site ID and location key
(?eipl,C02,000000,sid=AWS_ROOF&lkey=SITE01?)

# Set snap interval to 5 min and upload period to 15 min
(?eipl,C02,000000,mpm=5&upm=15?)

# Set primary server IP
(?eipl,C02,000000,sip1=192.168.1.100?)

# Set RTC to 09 March 2026, 10:30
(?eipl,C02,000000,rtc=2603091030?)

# Read all current system config values
(?eipl,C02,000000,?)
\end{verbatim}

\subsubsection{8.3 Sensor Enable / Disable
(C03)}\label{sensor-enable-disable-c03}

C03 enables or disables individual sensors and sets the wind sensor
protocol. Disabled sensors are excluded from the upload packet and CSV
log.

\textbf{Packet format:} \texttt{sdiid=enabled,cfgbits{[},uc{]}} groups
separated by \texttt{;}

\begin{longtable}[]{@{}lll@{}}
\toprule\noalign{}
sdiid & Sensor & Upload channels controlled \\
\midrule\noalign{}
\endhead
\bottomrule\noalign{}
\endlastfoot
\texttt{B} & Base ADC (ADS112C04 raw) & --- \\
\texttt{R} & Rain gauge & ch13 \\
\texttt{W} & Wind speed / direction (RS485) & ch9--ch12 \\
\texttt{T} & Temperature + Humidity (THP) & ch1--ch6 \\
\texttt{H} & Humidity / Pressure ADC slot & ch8 \\
\texttt{P} & Pressure / Solar ADC slot & ch7 \\
\texttt{S} & Solar radiation ADC slot & --- \\
\end{longtable}

\textbf{\texttt{cfgbits} for Wind sensor (\texttt{W}):}

\begin{longtable}[]{@{}lll@{}}
\toprule\noalign{}
cfgbits & Protocol & Sensor \\
\midrule\noalign{}
\endhead
\bottomrule\noalign{}
\endlastfoot
\texttt{0x00} & Modbus RTU & RENKEE \\
\texttt{0x01} & ASCII (Gill) & Gill Instruments \\
\end{longtable}

\textbf{\texttt{uc} (unit conversion) field:}

\begin{longtable}[]{@{}lll@{}}
\toprule\noalign{}
Sensor & uc value & Unit \\
\midrule\noalign{}
\endhead
\bottomrule\noalign{}
\endlastfoot
T (temperature) & 0 & °C (default) \\
T (temperature) & 1 & °F \\
W (wind speed) & 0 & m/s (default) \\
W (wind speed) & 1 & km/h \\
W (wind speed) & 2 & knots \\
W (wind speed) & 3 & mph \\
H (pressure slot) & 0 & Pa (default) \\
H (pressure slot) & 1 & hPa/mbar \\
H (pressure slot) & 2 & inHg \\
H (pressure slot) & 3 & mmHg \\
P (solar slot) & 0 & W/m² (default) \\
P (solar slot) & 1 & kWh/m² \\
\end{longtable}

\textbf{Examples:}

\begin{verbatim}
# Enable all sensors, RENKEE wind protocol, SI units
(?eipl,C03,000000,R=Y,0x00;W=Y,0x00;T=Y,0x00;H=Y,0x00;P=Y,0x00?)

# Enable Gill wind sensor
(?eipl,C03,000000,W=Y,0x01?)

# Disable rain sensor
(?eipl,C03,000000,R=N,0x00?)

# Set temperature to Fahrenheit (uc=1)
(?eipl,C03,000000,T=Y,0x00,1?)

# Set wind speed to km/h (uc=1)
(?eipl,C03,000000,W=Y,0x00,1?)

# Read all sensor config states
(?eipl,C03,000000,c03?)
\end{verbatim}

\subsubsection{8.4 Read Sensor Description
(C04)}\label{read-sensor-description-c04}

C04 is \textbf{read-only}. It returns the sensor slot configuration
(IDs, ranges, parameters) for diagnostic use:

\begin{verbatim}
(?eipl,C04,000000,?)
\end{verbatim}

\subsubsection{8.5 FTP Configuration (C05)}\label{ftp-configuration-c05}

C05 configures FTP upload credentials (reserved for future use):

\begin{longtable}[]{@{}ll@{}}
\toprule\noalign{}
Parameter & Description \\
\midrule\noalign{}
\endhead
\bottomrule\noalign{}
\endlastfoot
\texttt{ftu} & FTP username \\
\texttt{ftp} & FTP password \\
\texttt{ftd} & FTP directory \\
\texttt{fts} & FTP state \\
\end{longtable}

\begin{verbatim}
(?eipl,C05,000000,ftu=myuser&ftp=mypass&ftd=/data?)
\end{verbatim}

\subsubsection{8.6 Trigger Firmware Update
(C06)}\label{trigger-firmware-update-c06}

C06 initiates a FOTA firmware update check (same as the
\texttt{startfota} terminal command):

\begin{verbatim}
(?eipl,C06,000000,?)
\end{verbatim}

See \hyperref[12-firmware-over-the-air-fota-update]{Section 12} for full
FOTA details.

\subsubsection{8.7 THP Sensor Offsets
(C09)}\label{thp-sensor-offsets-c09}

Apply field correction offsets to the secondary
temperature/humidity/pressure sensor readings:

\begin{verbatim}
(?eipl,C09,000000,TEMP_offset=+0.5&HUM_offset=-2.0&PRESURE_offset=+1.5?)
\end{verbatim}

\begin{center}\rule{0.5\linewidth}{0.5pt}\end{center}

\subsection{9. Communication Interfaces}\label{communication-interfaces}

\subsubsection{9.1 USB CDC (Virtual COM
Port)}\label{usb-cdc-virtual-com-port}

The unit appears as a \textbf{USB CDC ACM} virtual serial port when
connected to a PC.

\begin{longtable}[]{@{}ll@{}}
\toprule\noalign{}
Property & Value \\
\midrule\noalign{}
\endhead
\bottomrule\noalign{}
\endlastfoot
Driver & Native (Windows 10/11, Linux, macOS) \\
Baud Rate & Not applicable (USB CDC) \\
Windows port name & \texttt{COMx} (Device Manager → Ports) \\
Linux/macOS port & \texttt{/dev/ttyACM0} or
\texttt{/dev/cu.usbmodem*} \\
Packet size & Up to 512 bytes \\
\end{longtable}

\textbf{Recommended terminal software:} PuTTY, TeraTerm, minicom, or any
serial terminal at any baud rate.

\textbf{Connection steps:} 1. Connect USB cable from PC to unit 2. Open
Device Manager; note the COM port assigned (e.g., COM8) 3. Open terminal
software, connect to COM8 (baud rate: any) 4. You will see the device
boot log or the \texttt{\textgreater{}} prompt 5. Send commands as
described in \hyperref[10-terminal-command-reference]{Section 10}

\subsubsection{9.2 UART1 Terminal}\label{uart1-terminal}

An alternate terminal interface is available on UART1 at \textbf{9600
baud, 8N1} (shared with the Quectel modem). This is primarily for
factory use or when USB is unavailable.

\subsubsection{9.3 RS485 Wind Sensor
Interface}\label{rs485-wind-sensor-interface}

The unit is the \textbf{Modbus master} on the RS485 bus. It queries the
wind sensor periodically. It does \textbf{not} act as a Modbus slave ---
external SCADA systems cannot query the ENVIRO\_DL over RS485.

\begin{longtable}[]{@{}llll@{}}
\toprule\noalign{}
Protocol & Baud & Format & Configuration \\
\midrule\noalign{}
\endhead
\bottomrule\noalign{}
\endlastfoot
RENKEE Modbus RTU & 9600 & 8N1 & \texttt{W=Y,0x00} \\
Gill ASCII & Configurable &
\texttt{"w?\textbackslash{}r\textbackslash{}n"} & \texttt{W=Y,0x01} \\
\end{longtable}

\begin{center}\rule{0.5\linewidth}{0.5pt}\end{center}

\subsection{10. Terminal Command
Reference}\label{terminal-command-reference}

\subsubsection{10.1 Command Packet Format}\label{command-packet-format}

All configuration commands use the Envirotech packet format:

\begin{verbatim}
(?eipl,<COMMAND_CODE>,<SEQUENCE>,<PAYLOAD>?)
\end{verbatim}

\begin{longtable}[]{@{}lll@{}}
\toprule\noalign{}
Field & Description & Example \\
\midrule\noalign{}
\endhead
\bottomrule\noalign{}
\endlastfoot
\texttt{eipl} & Fixed header & \texttt{eipl} \\
Command Code & C01--C24 & \texttt{C01} \\
Sequence & 6-digit counter (use \texttt{000000}) & \texttt{000000} \\
Payload & Key=value pairs delimited by \texttt{\&} &
\texttt{key=value\&key2=value2} \\
\end{longtable}

\subsubsection{10.2 Command Summary}\label{command-summary}

\begin{longtable}[]{@{}
  >{\raggedright\arraybackslash}p{(\linewidth - 4\tabcolsep) * \real{0.3000}}
  >{\raggedright\arraybackslash}p{(\linewidth - 4\tabcolsep) * \real{0.3333}}
  >{\raggedright\arraybackslash}p{(\linewidth - 4\tabcolsep) * \real{0.3667}}@{}}
\toprule\noalign{}
\begin{minipage}[b]{\linewidth}\raggedright
Command
\end{minipage} & \begin{minipage}[b]{\linewidth}\raggedright
Function
\end{minipage} & \begin{minipage}[b]{\linewidth}\raggedright
Direction
\end{minipage} \\
\midrule\noalign{}
\endhead
\bottomrule\noalign{}
\endlastfoot
\textbf{C01} & Station info (SID, time, FW version, signal, faults) &
Read \\
\textbf{C02} & System config: sid, lkey, mpm, upm, sip1, sip2, rtc, etc.
& Read/Write \\
\textbf{C03} & Sensor enable/disable, wind protocol (\texttt{cfgbits}),
unit conversion (\texttt{uc}) & Read/Write \\
\textbf{C04} & Sensor slot description (ranges, parameters) & Read \\
\textbf{C05} & FTP credentials (ftu, ftp, ftd, fts) & Read/Write \\
\textbf{C06} & Trigger FOTA firmware update check & Action \\
\textbf{C09} & THP sensor offsets (TEMP\_offset, HUM\_offset,
PRESURE\_offset) & Read/Write \\
\textbf{C10} & Current status report (latest sensor readings, all
channels) & Read \\
\textbf{C11} & Retrieve single day's log records from SD card & Read \\
\textbf{C12} & Retrieve multi-day log records by date range & Read \\
\textbf{C13} & Squall event status report & Read \\
\textbf{C21} & System reset & Action \\
\textbf{C22} & Clear SD card log & Action \\
\textbf{C24} & Reset rain counter (rain\_inday → 0) & Action \\
\textbf{help} & List all available commands & Read \\
\end{longtable}

\subsubsection{10.3 Short Terminal
Commands}\label{short-terminal-commands}

In addition to C-command packets, these short text commands are
accepted:

\begin{longtable}[]{@{}ll@{}}
\toprule\noalign{}
Command & Response \\
\midrule\noalign{}
\endhead
\bottomrule\noalign{}
\endlastfoot
\texttt{status} & Prints a full C10 status packet \\
\texttt{config} & Prints current configuration \\
\texttt{getclock} & Queries Quectel modem for network time \\
\texttt{getsim} & Returns SIM ICCID, IMSI, MSISDN \\
\texttt{getnet} & Returns network registration \& signal strength \\
\texttt{rawlog} & Enables raw UART1 hex capture to USB \\
\texttt{startfota} & Initiates a FOTA firmware update check \\
\end{longtable}

\subsubsection{10.4 C10 Status Response
Example}\label{c10-status-response-example}

\begin{verbatim}
(?eipl,R10A,000001,locationkey=SITE01&stationkey=AWS_ROOF&datetime=1741513500&ch1=26.50&ch2=25.30&ch3=24.10&ch4=65.20&ch5=64.50&ch6=63.80&ch7=450.00&ch8=1013.20&ch9=25.00&ch10=4.80&ch11=3.20&ch12=2.50&ch13=12.50&ch14=12.30?)
\end{verbatim}

\begin{longtable}[]{@{}lll@{}}
\toprule\noalign{}
Channel & Value in example & Meaning \\
\midrule\noalign{}
\endhead
\bottomrule\noalign{}
\endlastfoot
ch1 & 26.50 & Temperature maximum (°C) \\
ch2 & 25.30 & Temperature average (°C) \\
ch3 & 24.10 & Temperature minimum (°C) \\
ch4 & 65.20 & Humidity maximum (\%RH) \\
ch5 & 64.50 & Humidity average (\%RH) \\
ch6 & 63.80 & Humidity minimum (\%RH) \\
ch7 & 450.00 & Solar radiation (W/m²) \\
ch8 & 1013.20 & Barometric pressure (mbar) \\
ch9 & 25.00 & Wind direction (degrees) \\
ch10 & 4.80 & Wind speed maximum (m/s) \\
ch11 & 3.20 & Wind speed average (m/s) \\
ch12 & 2.50 & Wind speed minimum (m/s) \\
ch13 & 12.50 & Daily rainfall (mm) \\
ch14 & 12.30 & Battery/supply voltage (V) \\
\end{longtable}

\begin{center}\rule{0.5\linewidth}{0.5pt}\end{center}

\subsection{11. Remote Data Upload}\label{remote-data-upload}

\subsubsection{11.1 HTTP POST Upload}\label{http-post-upload}

The unit uploads sensor data to a remote server using an \textbf{HTTP
POST} request via the Quectel LTE modem.

\textbf{Request format:}

\begin{verbatim}
POST /api/data HTTP/1.1
Host: 192.168.1.100
Content-Type: application/x-www-form-urlencoded

datetime=20260309102500&locationkey=MUMBAI01&stationkey=AWS_ROOF&ch1=12.50&ch2=25.30...
\end{verbatim}

\textbf{Upload frequency:} Controlled by \texttt{upm} (upload period
minutes). The unit attempts upload at the end of each upload window. On
failure, it retries at the next interval.

\subsubsection{11.2 Upload Failure
Handling}\label{upload-failure-handling}

\begin{itemize}
\tightlist
\item
  If the server is unreachable, the record is queued and retry attempted
  next period
\item
  Successful uploads are marked in the log table to avoid duplicate
  sends
\item
  The upload counter on LCD Page 4 shows attempts vs.~successes for
  diagnostics
\end{itemize}

\subsubsection{11.3 LTE Modem Status}\label{lte-modem-status}

Use the terminal command \texttt{getnet} to check modem connectivity:

\begin{verbatim}
getnet
\end{verbatim}

Response includes registration status, operator name, and signal quality
(RSSI).

Use \texttt{getsim} to verify SIM is detected:

\begin{verbatim}
getsim
\end{verbatim}

\begin{center}\rule{0.5\linewidth}{0.5pt}\end{center}

\subsection{12. Firmware Over-the-Air (FOTA)
Update}\label{firmware-over-the-air-fota-update}

\subsubsection{12.1 Overview}\label{overview}

The ENVIRO\_DL supports firmware updates delivered by the Quectel LTE
modem from a remote FTP/HTTP server. The process is fully automated
with: - CRC32 integrity verification - Automatic rollback if the new
firmware is corrupted - Backup of the previous firmware before flashing

\subsubsection{12.2 Flash Memory Layout}\label{flash-memory-layout}

\begin{verbatim}
0x08000000  ┌─────────────────────────────────┐
            │        Bootloader  (128 KB)      │  Sectors 0–4
0x08020000  ├─────────────────────────────────┤
            │       Application  (896 KB)      │  Sectors 5–11
            │       ENVIRO_DL firmware         │
0x08100000  └─────────────────────────────────┘
\end{verbatim}

\subsubsection{12.3 Update Process}\label{update-process}

\textbf{Automatic (hourly check):} The \texttt{DataUploadManagerTask}
checks for a new firmware version every hour. If available, the update
begins automatically.

\textbf{Manual (terminal command):}

\begin{verbatim}
startfota
\end{verbatim}

\textbf{Stages:} 1. STM32 sends
\texttt{STARTFOTA:\textbackslash{}r\textbackslash{}n} to modem 2. Modem
responds
\texttt{FOTA:READY,\textless{}size\textgreater{},\textless{}version\textgreater{}\textbackslash{}r\textbackslash{}n}
(or \texttt{FOTA:NOTAVAIL\textbackslash{}r\textbackslash{}n} if no
update) 3. Firmware is streamed in 510-byte chunks to the SD card as
\texttt{update.bin} 4. CRC32 is verified against the value sent by the
modem 5. On CRC match: STM32 sets RTC backup flags and reboots 6.
Bootloader: - Backs up existing firmware to \texttt{backup.bin} on SD -
Erases application flash sectors - Programs and verifies
\texttt{update.bin} - On success: clears flags, jumps to new application
- On failure: restores \texttt{backup.bin}, continues with previous
firmware

\subsubsection{12.4 During an Update}\label{during-an-update}

\begin{itemize}
\tightlist
\item
  The LCD displays a \textbf{progress bar} on Page 5
\item
  The STATUS LED blinks rapidly
\item
  \textbf{DO NOT power off during a FOTA update}
\item
  If power is lost mid-update, the bootloader will detect incomplete
  flags and restore from \texttt{backup.bin} on the next power cycle
\end{itemize}

\subsubsection{12.5 Version Check}\label{version-check}

After a successful update, the firmware version is reported in the C10
status packet and in the boot log via USB.

\begin{center}\rule{0.5\linewidth}{0.5pt}\end{center}

\subsection{13. LED Status Indicator}\label{led-status-indicator}

The STATUS LED on GPIO PG4 provides visual indication of system health.

\begin{longtable}[]{@{}
  >{\raggedright\arraybackslash}p{(\linewidth - 2\tabcolsep) * \real{0.5000}}
  >{\raggedright\arraybackslash}p{(\linewidth - 2\tabcolsep) * \real{0.5000}}@{}}
\toprule\noalign{}
\begin{minipage}[b]{\linewidth}\raggedright
Pattern
\end{minipage} & \begin{minipage}[b]{\linewidth}\raggedright
Meaning
\end{minipage} \\
\midrule\noalign{}
\endhead
\bottomrule\noalign{}
\endlastfoot
\textbf{1 Hz blink} (500 ms on/off) & System healthy, normal
operation \\
\textbf{Solid ON} & Active fault (check LCD Page 0 or USB for fault
details) \\
\textbf{200 ms single flash} & Sensor reading completed \\
\textbf{3 × 100 ms rapid pulses} & Upload succeeded \\
\textbf{2-second solid} & Upload failed \\
\textbf{Rapid blink} & FOTA update in progress \\
\end{longtable}

\begin{center}\rule{0.5\linewidth}{0.5pt}\end{center}

\subsection{14. Fault Codes \&
Troubleshooting}\label{fault-codes-troubleshooting}

\subsubsection{14.1 Fault Bitmask}\label{fault-bitmask}

The system fault bitmask is reported in the C01 status packet
(\texttt{F=} field) and displayed on LCD Page 0. It is \textbf{not} a
numbered channel in the upload packet (ch13 is daily rainfall; ch14 is
battery voltage).

\begin{longtable}[]{@{}lll@{}}
\toprule\noalign{}
Bit (hex) & Fault & Cause \\
\midrule\noalign{}
\endhead
\bottomrule\noalign{}
\endlastfoot
\texttt{0x4000} & Power fault & Supply voltage out of range \\
\texttt{0x2000} & SD card fault & SD not mounted or write error \\
\texttt{0x1000} & Logging fault & File write failed \\
\texttt{0x0800} & RTC fault & Internal or external RTC unresponsive \\
\texttt{0x0400} & ADC fault & NAU7802 or ADS112C04 not responding \\
\texttt{0x0200} & Modem fault & Quectel modem not responding \\
\texttt{0x0100} & GPS fault & GPS query timeout (if GPS enabled) \\
\end{longtable}

A value of \texttt{0} means no faults. Read via:
\texttt{(?eipl,C01,000000,?)}

\subsubsection{14.2 Common Issues}\label{common-issues}

\textbf{Unit powers on but STATUS LED stays solid:} - Check USB terminal
for error messages - LCD Page 0 will show which subsystem
(\texttt{SD:ERR}, \texttt{MOD:ERR}, etc.) is faulted - Verify SD card is
inserted and formatted FAT32 - Verify sensor wiring (I2C pullups, RS485
polarity)

\textbf{No data on SD card:} - Confirm snap interval is configured
correctly (\texttt{C04}) - Check SD card is not full - Verify
\texttt{config.bin} is loaded at boot (USB terminal will log config load
status)

\textbf{Data not uploading:} - Check LTE signal: run \texttt{getnet} ---
should show ``registered'' and RSSI \textgreater{} -90 dBm - Verify
server IP and page configured correctly (\texttt{C06}) - Check SIM card
is active and data plan is enabled (\texttt{getsim}) - Review upload
counters on LCD Page 4

\textbf{Wind readings are zero:} - Verify RS485 wiring (A/B polarity) -
Confirm correct protocol selected (\texttt{W=Y,0x00} for RENKEE,
\texttt{W=Y,0x01} for Gill) - Check wind sensor power supply - Use USB
terminal --- boot log will show RS485 initialization status

\textbf{Rain counts stuck at zero:} - Verify rain gauge wire connected
to PG7 terminal - Test by momentarily shorting the rain gauge terminals
--- count should increment - Check rain gauge reed switch with
multimeter

\textbf{Time is wrong after power cycle:} - Set time via \texttt{C07}
command - Verify MCP7940 RTC has a coin cell battery installed - Check
USB boot log for \texttt{MCP7940\ sync} messages

\textbf{Temperature reading is out of range:} - Check PT1000 probe
wiring to NAU7802 ADC - Verify calibration gain/offset via C02 command -
Confirm NAU7802 I2C address (0x54) on bus

\begin{center}\rule{0.5\linewidth}{0.5pt}\end{center}

\subsection{15. SD Card File Reference}\label{sd-card-file-reference}

\begin{longtable}[]{@{}
  >{\raggedright\arraybackslash}p{(\linewidth - 2\tabcolsep) * \real{0.4348}}
  >{\raggedright\arraybackslash}p{(\linewidth - 2\tabcolsep) * \real{0.5652}}@{}}
\toprule\noalign{}
\begin{minipage}[b]{\linewidth}\raggedright
Filename
\end{minipage} & \begin{minipage}[b]{\linewidth}\raggedright
Description
\end{minipage} \\
\midrule\noalign{}
\endhead
\bottomrule\noalign{}
\endlastfoot
\texttt{DDMMYY.csv} & Daily log file (e.g., \texttt{090326.csv} = 9
March 2026) \\
\texttt{DDMMYY\_bkp.csv} & Backup of previous day's log \\
\texttt{config.bin} & Main configuration (snap interval, upload period,
server, sensor enables) \\
\texttt{nvdata.bin} & Non-volatile data (ADC calibration, battery
voltage cal, rain persistence) \\
\texttt{thp\_cfg.bin} & THP sensor offsets (temperature, humidity,
pressure corrections) \\
\texttt{ftp\_cfg.bin} & FTP server credentials (reserved for future
use) \\
\texttt{update.bin} & Staged FOTA firmware file (created during update,
deleted after flash) \\
\texttt{backup.bin} & Previous firmware backup (created by bootloader
before flashing) \\
\end{longtable}

\textbf{Maintenance tips:} - Periodically archive and delete old
\texttt{DDMMYY.csv} files to prevent the card from filling up - Do not
delete \texttt{config.bin}, \texttt{nvdata.bin}, or
\texttt{thp\_cfg.bin} --- these hold your calibration settings - Do not
remove the SD card while the unit is powered

\begin{center}\rule{0.5\linewidth}{0.5pt}\end{center}

\subsection{16. Technical
Specifications}\label{technical-specifications}

\subsubsection{Microcontroller}\label{microcontroller}

\begin{longtable}[]{@{}ll@{}}
\toprule\noalign{}
Parameter & Value \\
\midrule\noalign{}
\endhead
\bottomrule\noalign{}
\endlastfoot
MCU & STM32F417 \\
Core & ARM Cortex-M4 \\
Clock Speed & 168 MHz \\
Flash & 1 MB \\
RAM & 196 KB \\
RTOS & FreeRTOS (CMSIS-RTOS v2) \\
\end{longtable}

\subsubsection{Communication}\label{communication}

\begin{longtable}[]{@{}lll@{}}
\toprule\noalign{}
Interface & Standard & Notes \\
\midrule\noalign{}
\endhead
\bottomrule\noalign{}
\endlastfoot
USB & USB 2.0 Full-Speed CDC & Virtual COM port \\
LTE Modem & Quectel (AT command interface) & HTTP POST, FOTA \\
RS485 & Modbus RTU / ASCII & Wind sensor only \\
I2C & 100/400 kHz & Sensors, RTC \\
SPI & --- & LCD, SD card \\
\end{longtable}

\subsubsection{Sensors}\label{sensors}

\begin{longtable}[]{@{}llll@{}}
\toprule\noalign{}
Parameter & Sensor & Resolution & Range \\
\midrule\noalign{}
\endhead
\bottomrule\noalign{}
\endlastfoot
Temperature & PT1000 + NAU7802 & 24-bit ADC & -40 to +80 °C \\
Humidity & ADS112C04 & 24-bit ADC & 0--100 \%RH \\
Pressure & ADS112C04 / BMP280 & 24-bit ADC & 300--1100 hPa \\
Wind Speed & RS485 sensor & 0.1 m/s & 0--60 m/s (sensor dependent) \\
Wind Direction & ADS112C04 & 1° & 0--360° \\
Solar Radiation & ADS112C04 & --- & 0--2000 W/m² \\
Rainfall & Tipping bucket & 0.5 mm/tip & Unlimited accumulation \\
\end{longtable}

\subsubsection{Power}\label{power}

\begin{longtable}[]{@{}ll@{}}
\toprule\noalign{}
Parameter & Value \\
\midrule\noalign{}
\endhead
\bottomrule\noalign{}
\endlastfoot
Input Voltage & 5--12 V DC \\
Operating Current & \textasciitilde200--400 mA (typical, modem
active) \\
Sleep Current & N/A (always-on operation) \\
RTC Backup & External coin cell via MCP7940 \\
\end{longtable}

\begin{center}\rule{0.5\linewidth}{0.5pt}\end{center}

\subsection{17. Known Limitations}\label{known-limitations}

The following features are planned or partially implemented but not yet
available in this firmware version:

\begin{longtable}[]{@{}
  >{\raggedright\arraybackslash}p{(\linewidth - 4\tabcolsep) * \real{0.3750}}
  >{\raggedright\arraybackslash}p{(\linewidth - 4\tabcolsep) * \real{0.3333}}
  >{\raggedright\arraybackslash}p{(\linewidth - 4\tabcolsep) * \real{0.2917}}@{}}
\toprule\noalign{}
\begin{minipage}[b]{\linewidth}\raggedright
Feature
\end{minipage} & \begin{minipage}[b]{\linewidth}\raggedright
Status
\end{minipage} & \begin{minipage}[b]{\linewidth}\raggedright
Notes
\end{minipage} \\
\midrule\noalign{}
\endhead
\bottomrule\noalign{}
\endlastfoot
GPS / GNSS coordinates & Not implemented & Modem hardware capable;
firmware support pending \\
FTP upload & Not implemented & Infrastructure exists; not wired to
scheduler \\
Alert SMS on threshold breach & Not implemented & No real-time threshold
monitoring \\
Power-on SMS notification & Not implemented & --- \\
Network time sync (NTP) & Not implemented & RTC may drift over weeks;
set manually via C07 \\
Low-power / sleep mode & Not implemented & Unit is always-on; not
suitable for battery-only deployment without external management \\
USB pen drive data export & Not implemented & No USB MSC host stack \\
Real-time streaming (\texttt{stream\ N}) & Not implemented & Terminal
command not yet functional \\
Modbus slave (SCADA query) & Not implemented & Unit cannot be queried by
external RS485 SCADA \\
\end{longtable}

\begin{center}\rule{0.5\linewidth}{0.5pt}\end{center}

\subsection{Appendix A --- Quick-Start Command
Sequence}\label{appendix-a-quick-start-command-sequence}

\begin{verbatim}
# 1. Set station ID and location key (saved automatically)
(?eipl,C02,000000,sid=AWS_ROOF&lkey=SITE01?)

# 2. Set primary server IP
(?eipl,C02,000000,sip1=192.168.1.100?)

# 3. Set current date and time (format: YYMMDDHHmm)
(?eipl,C02,000000,rtc=2603091030?)

# 4. Set 5-minute snap interval and upload period
(?eipl,C02,000000,mpm=5&upm=5?)

# 5. Enable all sensors, RENKEE wind protocol, SI units
(?eipl,C03,000000,R=Y,0x00;W=Y,0x00;T=Y,0x00;H=Y,0x00;P=Y,0x00?)

# 6. Verify station info
(?eipl,C01,000000,?)

# 7. Read current sensor status
(?eipl,C10,000000,?)
\end{verbatim}

\begin{center}\rule{0.5\linewidth}{0.5pt}\end{center}

\subsection{Appendix B --- Glossary}\label{appendix-b-glossary}

\begin{longtable}[]{@{}
  >{\raggedright\arraybackslash}p{(\linewidth - 2\tabcolsep) * \real{0.3333}}
  >{\raggedright\arraybackslash}p{(\linewidth - 2\tabcolsep) * \real{0.6667}}@{}}
\toprule\noalign{}
\begin{minipage}[b]{\linewidth}\raggedright
Term
\end{minipage} & \begin{minipage}[b]{\linewidth}\raggedright
Definition
\end{minipage} \\
\midrule\noalign{}
\endhead
\bottomrule\noalign{}
\endlastfoot
\textbf{CDC} & USB Communications Device Class --- provides virtual COM
port \\
\textbf{FOTA} & Firmware Over-the-Air --- remote firmware update via
modem \\
\textbf{MMA} & Min/Max/Average --- statistical aggregation per snap
interval \\
\textbf{Modbus RTU} & Serial communication protocol used for wind
sensor \\
\textbf{PT1000} & Platinum resistance temperature detector, 1000 Ω at 0
°C \\
\textbf{RTC} & Real-Time Clock --- maintains date and time \\
\textbf{Snap interval} & Period between log file writes and data
uploads \\
\textbf{SHT25/SHT45} & Sensirion humidity + temperature sensor
(secondary) \\
\textbf{BMP280} & Bosch barometric pressure sensor (secondary) \\
\textbf{NAU7802} & 24-bit I2C ADC used for PT1000 temperature
measurement \\
\textbf{ADS112C04} & Texas Instruments 24-bit multi-channel ADC \\
\textbf{Squall} & Sudden wind speed increase \textgreater25\%, detected
by WMO state machine \\
\textbf{ch13} & Daily rainfall total (mm) in every upload packet \\
\textbf{ch14} & Battery / supply voltage (V) --- always included in
every upload packet \\
\end{longtable}

\begin{center}\rule{0.5\linewidth}{0.5pt}\end{center}

\emph{Document generated from firmware source code ---
ENVIRO\_DL\_STM32F, March 2026} \emph{For technical support, contact
SVIOT.}

\end{document}
